% =====================================================================
% chapters/06_discussion.tex
% Chapter 6 — Discussion
% Project: Vision-Based Running Technique Classification from In-the-Wild
%          Internet Videos using Markerless Pose Estimation and Deep Learning
%
% This chapter interprets the results reported in Chapter 5; it does not
% re-tabulate them. Numerical values are quoted only where needed to anchor
% an interpretation, and refer back to the controlled tables of Chapter 5.
% Controlled benchmark results and historical development results are kept
% strictly separate. No claim of real-world deployment readiness is made.
% =====================================================================

\chapter{Discussion}
\label{ch:discussion}

This chapter interprets the empirical results of Chapter~\ref{ch:results}. The
guiding question is not only \emph{how well} the pipelines perform, but
\emph{where} performance is gained and lost, and \emph{what that implies} for
the dataset and protocol used in this work.

% ---------------------------------------------------------------------
\section{Interpretation of the Main Findings}
\label{sec:disc-main}

The central pattern across all experiments is a decoupling between
\emph{representation learnability} and \emph{automatic segmentation
reliability}. When running cycles are correctly segmented, the markerless pose
representation is clearly sufficient: the handmade upper bound separates
correct from false technique at a macro-F1 near 0.90
(Section~\ref{sec:handmade}) . When the same representation and classifiers are
driven by automatically detected cycles, performance falls to roughly chance
(Section~\ref{sec:autocorr-baseline}) . The natural reading is that the
classification \emph{model} is not the binding constraint; the binding
constraint is the upstream pipeline that converts raw, uncontrolled video into
clean, comparable cycle samples.

This conclusion is supported at two levels. The fixed-split benchmark shows a
drop from 0.8991 macro-F1 for handmade cycles to 0.5417 for the autocorr
baseline. The fair-gap robustness check then repeats the comparison with the
same MLP classifier under 5-fold video-level cross-validation: handmade cycles
reach $0.860 \pm 0.016$ macro-F1, while autocorr cycles reach
$0.615 \pm 0.081$, leaving a fair gap of 0.245. Thus, the loss is not only a
single-split artefact; it remains visible when the comparison is repeated across
video-level folds.

This interpretation reframes the contribution of the thesis. The work does not
claim a high-accuracy automatic classifier. Instead, it provides evidence,
through a matched upper-bound/baseline design and a series of ablations, that
the difficulty of in-the-wild running-technique analysis lies predominantly in
robust cycle detection and pose quality rather than in the choice of classifier
or representation depth. Every subsequent section in this chapter can be read
as locating one facet of that bottleneck.

% ---------------------------------------------------------------------
\section{Effect of Cycle Source}
\label{sec:disc-cycle-source}

The comparison between the handmade and autocorrelation (autocorr) pipelines
(Section~\ref{sec:hm-vs-ac}) is the most informative single contrast in the
study. Because the two pipelines share the landmark representation, classifier
family, split, and reporting protocol, the large gap between them isolates the
effect of cycle segmentation almost cleanly: the manual and automatic pipelines
differ chiefly in \emph{how cycle boundaries are obtained}. The fair-gap
robustness check strengthens this interpretation: under the same MLP classifier
and 5-fold video-level cross-validation, handmade cycles achieve
$0.860 \pm 0.016$ macro-F1, whereas autocorr cycles achieve
$0.615 \pm 0.081$. The 0.245 macro-F1 fair gap shows that the fixed-split
drop is not simply due to one unlucky train/test partition.

Two mechanisms plausibly explain the gap. First, automatic detection discards a
substantial fraction of the data: many videos yield no usable cycle at all, and
of those that do, a large share of raw cycles are rejected by the
duration/quality filters (Section~\ref{sec:signal-ablation}). The surviving
cycles are therefore both fewer and biased toward easier cases. Second, even
retained cycles may have imprecise foot-strike-to-foot-strike boundaries, so
that the eight-phase samples are not temporally aligned across examples. A
classifier that performed well on manually aligned phases then faces a
distribution of misaligned, partially mis-segmented inputs. 

The validation-to-test collapse of the sequential models under automatic cycles
(Section~\ref{sec:autocorr-baseline}) is consistent with this misalignment: the
models latch onto regularities that do not transfer once boundaries vary. The
practical message is that improving the automatic pipeline should target
detection coverage and boundary precision before any change to the classifier.

The permutation test further tempers the interpretation of the autocorr signal.
The observed autocorr MLP CV macro-F1 was 0.5206, compared with a permutation
null mean of 0.468 and $p = 0.2475$. This does not prove that the automatic
features contain no information; rather, it suggests that the information is
weak and unstable under a simple MLP configuration. This is consistent with a
pipeline whose errors originate upstream in cycle boundary quality.

% ---------------------------------------------------------------------
\section{Effect of Representation: Landmark versus Image}
\label{sec:disc-representation}

On a matched automatic cycle source, the landmark representation clearly
outperformed the RGB-crop image representation
(Section~\ref{sec:representation}), and the image pipeline collapsed below
chance in the controlled rerun (Section~\ref{sec:image-results}). At first
sight this is surprising, since RGB crops carry strictly more raw information
(silhouette, limb contour, clothing-relative motion) than sparse landmarks.

The likely explanation is not that this information is useless, but that the
image pipeline cannot exploit it reliably under the present conditions: a small
training set, a high-capacity visual backbone prone to overfitting, and crop
instability driven by the same noisy pose estimates that bound the boundary
quality. In effect, the image pipeline inherits the segmentation bottleneck and
adds a second source of variance --- the runner-centred crop --- on top of it.

The landmark representation, by contrast, acts as a compact inductive bias:
by reducing each frame to a fixed set of joint coordinates, it discards
appearance variation that would otherwise have to be learned away from limited
data. For this dataset, that compression is an advantage rather than a loss.
The appropriate conclusion is conditional: image-based modelling is not
invalidated in principle, but the current implementation is not yet stable
enough to compete with the landmark baseline. A fair re-test would require more
data and more robust cropping.

This is also the right place to interpret the historical
\texttt{autocorr\_fixed} (val\_acc 76.5) milestone
(Section~\ref{sec:version-history}) . It represents the most technically ambitious
step in the project's development --- the move from landmark-only sequences to
an image-based MobileNetV2--LSTM pipeline --- and it motivated the organised
image-based rerun. However, its 76.5\% figure is a historical validation
accuracy from a development notebook, recorded under a different protocol and
never re-confirmed on the fixed test split. The controlled rerun did not
reproduce it. The milestone is therefore valuable as evidence of methodological
exploration and as the reason the image direction was worth testing, but it
cannot serve as a benchmark and must not be compared directly against the
controlled landmark baseline.

% ---------------------------------------------------------------------
\section{Effect of Model Architecture}
\label{sec:disc-model}

The model ablation (Section~\ref{sec:model-ablation}) shows that architecture
matters far less than data quality. On the clean handmade data, the CNN-BiLSTM
achieved the strongest macro-F1, while the 1D-CNN, BiLSTM, and MLP remained
close enough to show that the spread across architectures was narrow. This
indicates that, given well-segmented cycles, the discriminative signal is
accessible to several standard architectures and does not depend solely on deep
temporal machinery.

On the automatic pipeline the more expressive sequential models behaved
unstably. The 1D-CNN and BiLSTM both produced zero F1 for the false class at
their operating thresholds, while the CNN-BiLSTM looked best at the default
0.5 threshold but became unstable under validation-tuned thresholding. High-capacity sequence models
are precisely those most able to overfit the spurious regularities introduced
by inconsistent segmentation, so their poor generalisation is symptomatic of
the upstream noise rather than of an architectural deficiency. The robustness
of the simple MLP as the reported automatic baseline follows from the same
logic: with noisy, scarce inputs, a lower-variance model is preferable. This
argues against reaching for larger classifiers as a remedy for the automatic
pipeline's weakness.

% ---------------------------------------------------------------------
\section{Effect of Signal Choice}
\label{sec:disc-signal}

The signal ablation (Section~\ref{sec:signal-ablation}) exposed a tension
between two notions of a ``good'' detection signal. The vertical-oscillation
signals (\texttt{head\_y}, \texttt{hip\_y}) detected cycles in more videos and,
for \texttt{hip\_y}, even produced the highest downstream classification
macro-F1. Yet they admitted many ill-formed cycles, reflected in weak mean
autocorrelation and a large number of duration-based rejections.

The \texttt{contact\_diff} signal showed a different profile: fewer detected
videos, but cycle boundaries that were more directly tied to foot-contact
events and therefore more suitable for foot-strike-to-foot-strike sampling.
Selecting \texttt{contact\_diff} as the baseline signal therefore reflects a
deliberate preference for boundary \emph{precision} over detection
\emph{coverage}. This is the defensible choice for a pipeline whose downstream
representation is an eight-phase sampling that assumes accurate
foot-strike-to-foot-strike boundaries: a cleanly segmented minority is more
useful than a larger but temporally inconsistent set. 

The fact that \texttt{hip\_y} scored a higher macro-F1 is interpreted cautiously, given the
small test set and the risk that more permissive detection coincided with
easier videos rather than with genuinely better segmentation. The broader
lesson is that signal selection for in-the-wild data must be judged on
segmentation quality, not only on a single end-to-end accuracy number.

% ---------------------------------------------------------------------
\section{Effect of Feature Representation}
\label{sec:disc-feature}

The feature ablation (Section~\ref{sec:feature-ablation}) found that the
selected 13-landmark representation and the selected-landmark-plus-velocity
variant matched the best handmade macro-F1, while the full 33-landmark and
lower-body-only variants remained close but did not improve the overall
macro-F1. The conservative interpretation is that additional landmarks may
carry complementary cues for one class, but they did not produce a more robust
overall result on this split.

The near-equivalence of the selected-landmark
configurations supports using the compact representation as the pipeline
default, especially for the automatic experiments where every additional input
dimension also imports additional pose-estimation noise. This connects back to the representation discussion
(Section~\ref{sec:disc-representation}): both results favour parsimony.
Reducing the input to a small, well-chosen set of joints limits the model's
exposure to noise and to spurious appearance variation, which is advantageous
precisely because the dataset is small and the upstream estimates are imperfect.
The choice of features should therefore be coupled to the data regime rather
than maximised in isolation.

% ---------------------------------------------------------------------
\section{Effect of Frame Resolution per Cycle}
\label{sec:disc-frames}

The frames-per-cycle ablation (Section~\ref{sec:frames-ablation}) showed no
benefit from increasing temporal resolution beyond the eight-phase
representation, and there was no monotonic trend across frame counts. The
eight-phase sampling remained the strongest and most compact setting. The
interpretation is that eight phases already capture the salient kinematic
structure of a running cycle for this binary task; finer sampling mostly adds
redundant frames and, on limited data, additional opportunity to overfit rather
than additional discriminative signal.

This finding has a useful methodological consequence: the eight-phase
representation is justified not merely by convention but by an explicit
cost-benefit comparison. It keeps sample dimensionality low, which is
consistent with the parsimony favoured by the model and feature analyses, and
it reduces the computational burden of the pipeline without a meaningful
accuracy penalty.

% ---------------------------------------------------------------------
\section{Cycle-Level versus Video-Level Interpretation}
\label{sec:disc-cycle-video}

Aggregating cycle predictions to the video level by majority vote consistently
improved the automatic baseline (Section~\ref{sec:cycle-vs-video}), and the same
direction held for the other pipelines. This is the expected behaviour when
per-cycle predictions contain independent noise: aggregating several cycles from
the same video lets correct predictions outweigh sporadic errors, reducing
variance. Because the underlying data unit is the video, not the cycle,
video-level reporting is also the more faithful measure of practical
performance and is therefore the appropriate headline metric for any future
deployment-oriented evaluation.

The interpretation must nonetheless be tempered by sample size. The video-level
evaluation rests on only a handful of test videos, several of which contribute a
single cycle. The improvement from aggregation is consistent and theoretically
motivated, but its precise magnitude is uncertain at this scale, and the result
should be read as directional evidence rather than as a calibrated estimate.

% ---------------------------------------------------------------------
\section{Effect of Threshold Selection}
\label{sec:disc-threshold}

The threshold analysis (Section~\ref{sec:threshold}) revealed an asymmetry
between the clean and the automatic pipelines. On the handmade data,
validation-selected thresholds were mostly neutral or beneficial. On the
automatic data, the same procedure left the robust MLP unchanged but sharply
degraded the sequential models, with the CNN-BiLSTM losing substantial macro-F1
and being flagged as threshold-overfit.

This is the resolution of the apparent
contradiction with the model ablation: the CNN-BiLSTM looks best at the default
threshold but is unstable once a validation-chosen threshold is applied to
test, whereas the MLP is stable, which is why it is the reported baseline.
The deeper interpretation is that decision-threshold tuning is only as reliable
as the validation set on which it is performed. With a small validation split
and noisy automatic cycles, the validation-optimal threshold need not transfer,
and tuning can transfer noise rather than signal. 

The methodological
implication, reflected in the reporting throughout Chapter~\ref{ch:results}, is
that any tuned result must state that the threshold was selected on validation
and must be accompanied by the default-threshold result, so that threshold
overfitting is visible rather than hidden.

% ---------------------------------------------------------------------
\section{Effect of Pose and Cycle Quality}
\label{sec:disc-quality}

The quality analysis (Section~\ref{sec:quality}) provides the most direct
evidence for the segmentation-bottleneck interpretation. Classification
accuracy on the automatic pipeline rose with higher heel visibility, lower
landmark missing-ratio, and stronger autocorrelation of the cycle signal. In
other words, the errors are not distributed uniformly across the test set; they
concentrate in cycles where pose estimation is unreliable or the periodic
structure is weak.

Together with the failed-video breakdown, which attributes
the majority of failures to the cycle-detection stage, this locates the
dominant error source firmly upstream of the classifier. This is the unifying thread of the discussion. The cycle-source gap
(Section~\ref{sec:disc-cycle-source}), the fragility of expressive models under
automatic cycles (Section~\ref{sec:disc-model}), the preference for the
boundary-precise \texttt{contact\_diff} signal
(Section~\ref{sec:disc-signal}), and the instability of the image pipeline
(Section~\ref{sec:disc-representation}) are all consequences of the same
underlying limitation: imperfect markerless pose estimation and the difficulty
of detecting clean running cycles in uncontrolled footage. The quality analysis
makes that limitation measurable.

% ---------------------------------------------------------------------
\section{Error Analysis}
\label{sec:disc-error}

Synthesising the error evidence, three failure modes can be distinguished. First, \emph{detection failures}: videos in which no usable cycle is found, so
the system cannot make a prediction at all; these dominate the failed-video
counts and arise mainly at the cycle-detection stage
(Section~\ref{sec:quality}) . Second, \emph{segmentation errors}: cycles that are
detected but poorly bounded, producing misaligned eight-phase samples that
degrade classification even when pose visibility is adequate; the
validation-to-test collapse of the sequential models is the clearest symptom
(Section~\ref{sec:autocorr-baseline}) . Third, \emph{classifier errors} on
genuinely clean cycles, which the handmade results show to be comparatively
rare.

The relative weight of these modes is itself a result: the first two,
both upstream, account for the bulk of the gap between the upper bound and the
automatic baseline, while the third is small. This ordering should guide
remediation. Reducing detection and segmentation errors --- through more robust
pose estimation, better signal conditioning, or a learned boundary detector
trained on more data --- is likely to yield far larger gains than further work
on the classifier. The exploratory boundary-detector and raw-sequence pipelines
(Section~\ref{sec:bd-raw}) are first steps in that direction, but on the present
data they did not yet improve over the rule-based baseline and remain
inconclusive.

% ---------------------------------------------------------------------
\section{Limitations}
\label{sec:disc-limitations}

Several limitations bound the conclusions of this work. The dataset is small,
and the validation and test splits in particular are very small
(Section~\ref{sec:dataset-stats}) ; consequently, single-point differences,
figures after threshold tuning, and the highest-scoring feature configuration
should be treated as indicative rather than precise. The 5-fold fair-gap check
partly mitigates this concern by repeating the main handmade--autocorr contrast
across folds, but it does not remove the need for a larger external dataset. The video-level results
also rest on few videos. The class distribution is moderately imbalanced at the
video level, which is mitigated but not eliminated by the use of macro-F1
\cite{sokolova2009metrics}.

The pipeline also inherits the limitations of its components. Markerless pose
estimation is itself noisy on in-the-wild footage \cite{bazarevsky2020blazepose}.
Recent markerless-running analyses also note that such noise can affect
downstream biomechanical interpretation \cite{panconi2024markerless}. In this
pipeline, that noise propagates into signal extraction, cycle detection, and
--- for the image pipeline --- crop stability. 

Cycle detection relies on a
periodicity assumption that does not hold for very short clips or atypical gait,
contributing to detection failures. The handmade pipeline, used as the upper
bound, depends on manual segmentation and therefore reflects an idealised
condition that an automatic system cannot fully attain. Finally, all results
are obtained on a single corpus collected under one labelling convention;
generalisation to other populations, camera viewpoints, and definitions of
``correct'' technique remains untested. The historical version-history results
(Section~\ref{sec:version-history}) are excluded from all benchmark claims for
the reasons given above and add a further caution against over-reading
development-stage numbers. The permutation test is also limited by the same
small autocorr sample size and should be interpreted as a diagnostic robustness
check, not as a definitive statistical claim about all possible classifiers.

% ---------------------------------------------------------------------
\section{Practical Implications}
\label{sec:disc-practical}

The practical implications follow directly from the bottleneck analysis. For
researchers and practitioners building on this work, effort is best invested in
the upstream stages: more robust markerless pose estimation, signal conditioning
that improves cycle-boundary precision, and data-driven boundary detection
trained on enough examples to outperform rule-based autocorrelation. The
evidence indicates that such improvements would raise the automatic pipeline
toward the handmade upper bound more effectively than any change to the
classifier or any increase in temporal or feature resolution. Where an
automatic system is used, predictions should be reported and consumed at the
video level, with explicit handling of videos for which no reliable cycle can
be detected.

It must be stated plainly that these results do not establish deployment
readiness. The automatic pipeline performs near chance on the controlled test
set, the evaluation rests on a small single-source dataset, and the upper bound
depends on manual segmentation that a deployed system would not have. 

The
appropriate framing of the contribution is diagnostic: the work demonstrates
that markerless pose sequences can encode running-technique distinctions when
cycles are clean, and it identifies, through controlled comparison and
ablation, exactly where an automatic in-the-wild system currently breaks down.
Turning that diagnosis into a usable tool is left to future work
(Chapter~\ref{ch:conclusion}) .

% End of Chapter 6